\title{Position: There Are Futures That Benchmark-Driven AI Cannot See}
\begin{document}
\maketitle

\begin{abstract}
In biology, traits evolved for one function some- times become decisive for another; this is exapta- tion, and scientific progress works the same way. AI’s benchmark-centered selection environment taxes exaptation: when one selection rule dom- inates, ideas that do not fit it have nowhere to persist. The cost grows acute as the field shifts from asking can machines exhibit intelligent be- havior? to asking can they do so while remaining aligned, interpretable, and safe? These are philo- sophically distinct questions that may require dis- coveries we cannot specify. We propose mecha- nisms to restore exaptive capacity without aban- doning benchmarking: plural evaluation regimes, protected venues for non-comparable work, long- horizon funding, and training norms that encour- age researchers to question selection rules, not only optimize within them.
\end{abstract}

\section{Recovered Text}
Position: There Are Futures That Benchmark-Driven AI Cannot See
Sobhan Lotfi * 1 Ava Iranmanesh * 2 Lachin Naghashyar 3 Ali Shirali 4 Fateme Nateghi Haredasht 5
Sanmi Koyejo 5 Philip Torr 6 Yong Suk Lee 7 Fazl Barez 6 Joel Lehman 8 Peter Norvig 5 9 Arvind Narayanan 10
Abstract
In biology, traits evolved for one function some-
times become decisive for another; this is exapta-
tion, and scientific progress works the same way.
AI’s benchmark-centered selection environment
taxes exaptation: when one selection rule dom-
inates, ideas that do not fit it have nowhere to
persist. The cost grows acute as the field shifts
from asking can machines exhibit intelligent be-
havior? to asking can they do so while remaining
aligned, interpretable, and safe? These are philo-
sophically distinct questions that may require dis-
coveries we cannot specify. We propose mecha-
nisms to restore exaptive capacity without aban-
doning benchmarking: plural evaluation regimes,
protected venues for non-comparable work, long-
horizon funding, and training norms that encour-
age researchers to question selection rules, not
only optimize within them.
1. Introduction
“Live long and be patient; this wheel of tricks has
a thousand wilder games to play.”
— Hafez Shirazi
Modern machine learning is largely organized around bench-
mark performance. This has proven to be an unusually
effective epistemic choice: shared tasks, datasets, and met-
rics have made progress directly comparable across labs,
legible to funders, and easier to coordinate as the field ex-
*Equal contribution
1Department of Computer Engineering,
Sharif University of Technology, Tehran, Iran 2Pennsylvania State
University, State College, PA, USA 3Microsoft, London, UK
4University of California, Berkeley, CA, USA 5Stanford University,
Stanford, CA, USA 6University of Oxford, Oxford, UK 7Keough
School of Global Affairs, University of Notre Dame, Notre Dame,
IN, USA 8Lila Sciences, San Francisco, CA, USA 9Google, Moun-
tain View, CA, USA 10Princeton University, Princeton, NJ, USA.
Correspondence to: Ava Iranmanesh <ami5520@psu.edu>.
Proceedings of the 43 rd International Conference on Machine
Learning, Seoul, South Korea. PMLR 306, 2026. Copyright 2026
by the author(s).
panded (Jelinek, 1997; Harman, 1993). Rather than resolv-
ing longstanding philosophical debates about the nature of
intelligence, benchmarking offered a pragmatic alternative:
measuring what systems can objectively do (Koch \& Pe-
terson, 2024). We do not argue that benchmarking was a
mistake; we ask what is lost when evaluation is concentrated
around a single selection rule. Any dominant criterion privi-
leges what it can measure while sidelining what it cannot;
the resulting omissions are easy to overlook because the
system otherwise appears to be working. As Kuhn noted,
periods of “normal science” routinely set aside anomalies;
those anomalies can later become central when a paradigm
enters crisis and shifts (Kuhn, 1962; Arthur, 1996).
Our lens is exaptation, a biological term for traits evolved for
one function that later become decisive for another (Gould
\& Vrba, 1982). In research ecosystems, ideas can similarly
acquire value in ways their originators did not anticipate
(Wimsatt, 2007; Stanley \& Lehman, 2015). Crucially, exap-
tation depends on persistence: ideas with latent potential
must survive periods when they appear uncompetitive by
prevailing metrics in order to later realize that potential. We
argue that benchmark-centered evaluation imposes an
exaptation tax by reducing the survival of work whose
value is latent, indirect, or only revealed under a future
problem regime. In this sense, benchmarks do not merely
evaluate research; they organize it, shaping who can partici-
pate, what kinds of contributions are rewarded, and which
questions and methods come to seem fundable, publishable,
and legitimate (Bourdieu, 1975).
The stakes of this tradeoff rise as the field’s guiding question
changes. For much of its history, the implicit target was Q1:
can machines exhibit intelligent behavior? Increasingly, the
focus has shifted toward Q2: can machines exhibit intelli-
gent behavior such that they are aligned, interpretable, safe,
and governable? These “such-that” clauses are not simply
additional constraints on the original question. They change
what counts as success: Q1 has ground truth that exists in
the world independently of who evaluates the system; Q2
does not. Whether a system is aligned, interpretable, or
safe depends on criteria that must be constructed through
deliberation rather than read off from the task itself (Russell
\& Norvig, 2021; Gabriel, 2020). The field has not yet estab-
lished that the correlation between benchmark performance
1
Futures that benchmark-driven AI cannot see
and real progress that held for Q1 holds for Q2. Importing
Q1’s evaluation framework into Q2 without validating that
it transfers is not a conservative choice; it forecloses the
exploratory work from which any adequate answer to Q2
would have to emerge.
The remainder of the paper traces how benchmark-centered
evaluation became dominant, showing that its current cen-
trality is historically contingent rather than inevitable (Sec-
tion 2); diagnoses the mechanisms by which it governs
research trajectories to clarify where intervention is possible
(Section 3); develops the such-that shift as a problem that
differs from Q1 not in degree but in kind (Section 4); and
proposes institutional mechanisms for restoring exaptive ca-
pacity without abandoning benchmarking (Section 5). The
cost of getting this wrong is not a missed benchmark; it
is a field that optimizes confidently toward futures it can
measure, and misses the ones it needs most.
Conflict of Interest Disclosure: Lachin Naghashyar is
employed by Microsoft. Joel Lehman is employed by Lila
Sciences. Peter Norvig is employed by Google. The authors
declare that this research was conducted in the absence
of any commercial or financial relationships that could be
construed as a potential conflict of interest.
2. Background: Exaptation and AI History
This section reviews exaptation and the history of AI evalu-
ation in order to show how selection regimes shape which
ideas persist long enough to matter.
2.1. Exaptation and the Persistence of Ideas
Exaptation was formally introduced by Gould and Vrba as
a refinement of adaptation, referring to traits that evolved
for one function and were later co-opted for another (Gould
\& Vrba, 1982). Exaptation depends on persistence beyond
a trait’s original function. Such traits may be adaptive in
their original role, but that fitness does not explain their
later significance. In early mammals, for example, bones
that once played a role in jaw mechanics were later repur-
posed into the middle ear, enabling hearing (Anthwal et al.,
2013). Their importance in this new role was not explained
by optimization for the original one, but by their survival
long enough to be reused. This logic has been extended be-
yond biology to scientific and technological change, where
ideas or tools developed for one problem can later become
decisive for another, provided they persist through periods
in which their value lies outside what prevailing evaluative
standards reward (Ferreira et al., 2020). Magnetrons de-
veloped for radar became the basis for microwave cooking.
Non-Euclidean geometry, developed by Riemann in 1854 as
pure mathematics with no physical application, became es-
sential to general relativity sixty years later. Exaptation thus
highlights a mode of progress in which long-term innovation
depends not only on optimization for current objectives, but
on the continued institutional viability of ideas whose value
is not yet legible.
2.2. Evaluation Regimes in the History of AI
In the pre-benchmark era of AI (1956–1987), the field de-
veloped as a pluralistic but loosely evaluated research space,
with competing theories of intelligence and no shared stan-
dards for evaluating progress (Russell \& Norvig, 2021). The
Dartmouth proposal articulated an expansive ambition, as-
serting that “every aspect of learning or any other feature
of intelligence can in principle be so precisely described
that a machine can be made to simulate it”, without speci-
fying how progress should be assessed across approaches
(McCarthy et al., 1955). This openness allowed symbolic,
connectionist, and heuristic methods to coexist, but left
no shared standard for resolving disputes. Early AI has
been characterized as a form of “protoscience,” in which or-
ganic evaluation mechanisms proved insufficient to stabilise
consensus, contributing to recurring hype–disappointment
cycles and the AI winters of the 1970s and late 1980s.
In the benchmark intervention period (1987–2012), major
funders—most notably DARPA—helped reorient AI away
from resolving foundational disputes about intelligence and
toward measurable progress on well-specified tasks with
clear evaluation metrics. This shift crystallised into the
Common Task Framework, which organized research around
shared datasets, held-out evaluation, and automated scor-
ing, enabling direct comparison among methods on a fixed
task. In this sense, benchmarking served as a pragmatic
bypass, making progress legible through quantitative perfor-
mance (Koch \& Peterson, 2024; Raji et al., 2021a).
From roughly 2012 onward, AI entered a period of extraordi-
nary coherence driven by the convergence of large datasets,
GPU acceleration, and accumulated advances in learning
algorithms. Following the 2012 ImageNet breakthrough,
deep learning’s success reflected favourable positional fac-
tors—large labelled datasets and GPU acceleration—rather
than a settled theory of intelligence (Krizhevsky et al., 2012).
Once this combination proved decisive across vision, speech,
and language tasks, scaling model size, data, and compute
emerged as a reliable and comparatively low-risk path to
performance gains. Empirical scaling studies reinforced
this strategy by demonstrating predictable improvements
with increased scale, further entrenching optimization over
a small set of quantitative benchmarks as the dominant re-
search practice (Hestness et al., 2017). As a result, the field
increasingly treated benchmark performance as the main
indicator of progress, yielding rapid gains and clear success
signals, and producing an evaluation regime tightly focused
on predictive accuracy. By mid-2026, leading benchmarks
2
Futures that benchmark-driven AI cannot see
have largely saturated, making this regime’s limits visible
rather than merely theoretical.
2.3. Exaptation in Modern Machine Learning
The historical record suggests that evaluation regimes play a
role in determining which lines of work remain visible, dis-
seminated, and sustained by the research community long
enough to matter. Modern machine learning provides con-
crete illustrations of this pattern: several of its core capabili-
ties emerged through the repurposing of ideas, tools, or struc-
tures, sometimes developed for adjacent problems or sub-
fields, and later applied in new roles. At the infrastructural
level, graphics processing units (GPUs) provide a canoni-
cal example. Originally developed to accelerate real-time
graphics rendering, GPUs were optimised for highly par-
allel arithmetic operations rather than for general-purpose
programmability (Owens et al., 2008; Nickolls et al., 2008).
This hardware proved decisive for deep learning by making
the training of previously impractical neural networks fea-
sible at scale. The 2012 ImageNet breakthrough, enabled
by GPU-trained convolutional networks, marked a turning
point that accelerated the adoption of deep learning across
vision, speech, and language (Krizhevsky et al., 2012). In
this sense, GPUs did not constitute a theoretical advance in
learning algorithms, but an infrastructural exaptation that re-
shaped the pace and direction of machine learning research.
Exaptation has also operated at the levels of model archi-
tecture and training methodology. The U-Net architecture
was originally introduced for biomedical image segmenta-
tion (Ronneberger et al., 2015). In later work on denoising
diffusion probabilistic models, U-Net–based architectures
were repurposed as the core backbone for generative model-
ing, where the same multi-scale, skip-connected structure
supports iterative noise prediction and removal, enabling
high-fidelity image synthesis (Ho et al., 2020).
A useful illustration comes from the history of neural net-
works in the late 1990s and early 2000s. Support vector
machines appeared stronger on the benchmarks of the era,
and the field shifted attention accordingly. This was locally
rational but ultimately misleading: advances in optimization,
data availability, and GPU acceleration later enabled deep
networks to become the foundation of modern ML (Bengio
et al., 2013; LeCun et al., 2015). The episode illustrates how
evaluation regimes shape which ideas are prematurely set
aside.
The 2014 paper introducing generative adversarial networks
offers a more recent and precisely documented case (Good-
fellow et al., 2014). The authors hedged explicitly on sample
quality, writing that their results made no claim of superi-
ority over existing generative methods, framing the con-
tribution as demonstrating a framework’s potential. One
reviewer concluded that the submission was “may not yet
be quite strong enough for NIPS.” The paper was accepted.
It became one of the most cited works in the history of ma-
chine learning. What the evaluation regime could not easily
measure was the generative framework itself; the conceptual
architecture that would later reshape the field.
3. Present Diagnosis: The Tyranny of
Benchmarks
Benchmarks do not merely evaluate research, they increas-
ingly organize it. This governance operates on three levels.
First, benchmark-centered progress selects for particular
resource profiles (compute, engineering infrastructure, fa-
miliarity with standardized datasets), shaping who can par-
ticipate and who can sustainably stay. Second, benchmarks
shape what even the selected participants work on, systemat-
ically rewarding “hill-climbable” contributions while raising
the cost of paradigm-challenging work. Third, benchmarks
codify present understanding, while many breakthrough con-
tributions derive their value from future exaptations, uses
and recombinations that simplistic evaluation cannot antic-
ipate, yet simplistic evaluation determines which lines of
work survive long enough to be exapted.
The tyranny is not in any individual benchmark. It is in
benchmark culture as the field’s dominant mode of eval-
uation: when leaderboards, standardized pipelines, and
common-task success become the default currency of le-
gitimacy, they quietly reshape the community, the agenda,
and the kind of novelty that can be recognized.
3.1. Community Composition
NeurIPS was founded by researchers studying neural in-
formation processing systems, a coalition spanning neuro-
science, cognitive science, physics, and computer science,
unified by interest in how biological and artificial systems
process information. That founding coalition has largely
moved to the periphery of today’s flagship ML venues. Bour-
dieu offers a framework for understanding such shifts: sci-
entific fields define what counts as legitimate capital, and
participants without the right capital learn to self-disqualify
before explicit rejection occurs (Bourdieu, 1975). In ML
today, that cost of entry is not only conceptual; it is often
infrastructural.
This narrowing is not incidental to benchmark culture—it is
produced by it. Klinger et al. document that AI research has
become focused on narrower themes in recent years, and
that work involving the private sector tends to be less di-
verse and more influential than academic work, with special-
ization toward data-hungry and computationally intensive
methods (Klinger et al., 2020; OECD, 2023). A large-scale
analysis of 41 million research papers confirms that AI tools
amplify individual output while collectively narrowing the
3
Futures that benchmark-driven AI cannot see
themes researchers pursue (Hao et al., 2026). This con-
centration also reflects a structural feature of modern AI’s
funding landscape: unlike earlier technology breakthroughs,
from the internet to GPS to semiconductors, which were
primarily government- and academia-funded (Mazzucato,
2013), today’s AI development is driven by venture capital
and large technology companies whose selection pressures
favor commercially viable, short-horizon research (OECD,
2023). Benchmark culture amplifies this dynamic by mak-
ing large-scale compute, data, and engineering infrastructure
especially valuable: when progress is measured as incre-
mental gains on shared tasks, the returns to scale in these
resources become unusually direct, and compatibility with
prevailing hardware and software ecosystems shapes which
ideas succeed as much as their scientific merit (Hooker,
2020; Strubell et al., 2019).
Why should we care who participates? One reason is com-
binatorial: transformative work often emerges from unusual
recombinations across methods, domains, and problem fram-
ings. The empirical science-of-science literature supports
this mechanism at scale. Wu, Wang, and Evans, analyzing
65 million papers and patents, find that small teams more of-
ten “disrupt” existing trajectories while large teams more of-
ten “develop” and consolidate them (Wu et al., 2019). Uzzi
et al. similarly show that the highest-impact science tends
to combine highly conventional prior work with a smaller
intrusion of atypical combinations (Uzzi et al., 2013). The
implication for ML is not that large-scale benchmarked en-
gineering is bad (it is often essential), but that a community
whose dominant selection pressures favor scale and consoli-
dation will, at the margin, generate fewer atypical combina-
tions, and atypical combinations are the raw material from
which future exaptations are built.
Consider what this means for exaptation in ML itself. The
attention mechanism emerged from work at the intersection
of machine learning and computational linguistics — a com-
munity that maintained distinct evaluation standards and in-
stitutional identity for decades (Bahdanau et al., 2014). ACL
remains one of AI’s oldest venues, and translation one of
its oldest problems. That community sustained the concep-
tual foundations, sequence modeling, alignment, encoder-
decoder structure, later extended into the transformer (Bah-
danau et al., 2014; Vaswani et al., 2017). AlphaFold2 then
adapted these architectures to protein structure prediction,
recognized by the 2024 Nobel Prize in Chemistry (Jumper
et al., 2021; The Nobel Prize, 2024). The lineage is coun-
terintuitive: solving a grand challenge in structural biology
required having had linguists in AI. This was not planned,
and could not have been. It was possible because a hetero-
geneous research ecosystem preserved communities with
distinct problem framings long enough for unexpected re-
combinations to occur. Benchmark culture can sustain nar-
row subcommunities. But cross-domain exaptations depend
on the broader ecosystem that benchmark monoculture se-
lects against.
This is why “community composition” is not a sociological
side note: it is a technical precondition for certain kinds of
progress. If benchmark culture disproportionately rewards
resource-intensive, pipeline-compatible work, then it indi-
rectly shapes the space of recombinations the field will ever
get a chance to discover.
3.2. Research Agenda
Section 3.1 established that benchmark culture shapes who
participates in ML research. But benchmarks also shape
what even selected participants choose to work on. Donoho
articulates a compelling positive case: “frictionless repro-
ducibility”, the combination of shared data, shared code,
and competitive challenges, can accelerate progress by en-
abling researchers to directly compare, iterate, and build
with minimal coordination cost (Donoho, 2024). This is
real and valuable. Yet Recht, in his commentary on Donoho,
offers a revealing reframe: “We can view machine learning
practice as a massively parallel genetic algorithm that fixates
on goals where hill climbing is possible and ignores those
where it isn’t” (Recht, 2024).
This dynamic maps naturally onto Kuhn’s distinction be-
tween normal science and revolutionary science (Kuhn,
1962).
Benchmarks, leaderboards, and standardized
pipelines are “normal science with the friction removed”:
they create extraordinary efficiency at paradigm-consistent
progress. But the same infrastructure raises the barrier to
work that challenges the paradigm. To depart from the
dominant framing, one must often do two things at once:
introduce a new idea and demonstrate competitive perfor-
mance on metrics and pipelines optimized for the incumbent
paradigm. The asymmetry is structural: mature baselines
inherit years of tuned hyperparameters, training recipes,
compute infrastructure, and evaluation conventions, while
genuinely new approaches must recreate much of that stack
merely to become comparable.
The transition from statistical to neural machine translation
shows this asymmetry playing out over an entire research
generation. Neural approaches existed early (Forcada \&
˜Neco, 1997), but SMT dominated for decades under bench-
mark regimes and engineering ecosystems that strongly fa-
vored it (Stahlberg, 2020). NMT “arrived” as a benchmark
winner only after the broader neural ecosystem (compute,
data, training methods, and engineering practices) matured
to the point where neural systems could become competitive
quickly (Sutskever et al., 2014; Bahdanau et al., 2014). The
lesson is not that benchmarks prevented NMT (they even-
tually validated it). The lesson is that benchmark-driven
selection pressures can delay or tax paradigm shifts until the
surrounding ecosystem makes them “one-shot competitive.”
4
Futures that benchmark-driven AI cannot see
Recent analyses confirm that benchmarks steer agendas,
not merely measure outcomes. Systematic study of sixty
LLM benchmarks finds that saturation is now the norm:
once models approach ceiling performance, the field’s re-
sponse is to create new benchmarks rather than to question
the evaluative regime itself (Akhtar et al., 2026; Ott et al.,
2022). Leaderboard rankings have been shown to reflect
undisclosed private testing practices that benefit a hand-
ful of providers, distorting the playing field that nominally
drives research priorities (Singh et al., 2025). Critiques
from within ML further emphasize construct validity prob-
lems when benchmarks are treated as broad measures of
progress (Raji et al., 2021b; Bechler-Speicher et al., 2025).
Benchmarks function as selection environments that shape
which questions are “worth asking” and which styles of
research are “viable careers.”
3.3. The Limits of Evaluation
Every evaluation instantiates a theory of what matters. That
theory is necessarily incomplete because the future problem
landscape is not knowable in advance. This is not merely or-
dinary uncertainty (unknown parameters); it is often closer
to Knightian uncertainty: qualitative shifts in tasks, environ-
ments, and use cases that cannot be exhaustively enumerated
today (Knight, 1921). Lehman et al. make this point in a
particularly ML-relevant form, arguing that ML systems
face a “Knightian blindspot” when confronted with qual-
itatively novel challenges generated by complex, creative
environments (Lehman et al., 2025). When evaluation is
optimized for what we can currently specify, it can system-
atically underweight work whose value depends on what we
cannot yet quantitatively specify.
Historical patterns in science and technology recur. Work
can be “ahead of its time,” recognized only when com-
plementary tools or applications arrive (van Raan, 2004;
Garfield, 1980). In AI, the most valuable contributions can
be those that later become indispensable components of sys-
tems aimed at problems we do not yet know how to write
down.
Benchmarks also encode hidden assumptions about deploy-
ment. Computer vision benchmarks, for example, histor-
ically defined tasks through particular label sets, sensors,
and controlled conditions—assumptions that later turned
out to be misaligned with robotics, distribution shift, and
real-world robustness (Everingham et al., 2010; Scharstein
\& Szeliski, 2002; Hendrycks \& Dietterich, 2019). This does
not indict benchmarking; it clarifies its limits: benchmarks
measure what they measure.
Path dependence theory formalizes why such limits can have
long-run consequences. When a research direction experi-
ences increasing returns (tooling, shared pipelines, network
effects in adoption), early advantages compound and alterna-
tives become harder to pursue even if they would be better in
retrospect (Arthur, 1989; 1996). Benchmark culture creates
increasing returns for specific research profiles: compute,
engineering infrastructure, and familiarity with standard
datasets. These advantages accumulate. The resulting loss
is largely invisible because we cannot observe innovations
never pursued. Yet the structural dilemma remains: bench-
marks codify present understanding; breakthroughs often
require violating that understanding. No quantitative eval-
uation method can foresee exaptations. But our evaluation
methods determine which research survives long enough to
be exapted.
4. The Such-That Problem
The three dimensions diagnosed in Section 3 (community
composition, research agenda, and the limits of evaluation)
each become more acute as the field’s guiding question
shifts from Q1 to Q2, and the worsening is not one of de-
gree. Community composition matters more because such-
that problems require contributions from disciplines that
benchmark culture systematically marginalizes: philosophy,
political science, and the social sciences that have grappled
with contested normative questions for decades. Research
agenda matters more because benchmark-driven selection
cannot orient toward goals that are not yet fully specified;
one cannot hill-climb toward a summit whose location re-
mains unknown. The limits of evaluation become the central
problem rather than a manageable one: for capability re-
search, the gap between benchmark performance and real
progress is imperfect but has been validated by decades of
use; for such-that problems, that validation does not exist,
and premature convergence on current proxies may fore-
close discovery of what we actually need to measure. The
following subsections develop why the shift from Q1 to Q2
changes the nature of the problem, not only its difficulty.
4.1. The Such-That Problem Is in the Air
We cite what follows as evidence that this shift is already
institutional, not as an evaluation of the efforts themselves.
The arc of this section is observational: 4.1 documents the
shift, 4.2 argues it is different in kind from what preceded it,
and 4.3 draws the implication for exaptation.
Evidence suggests a shift in research attention.
Tur-
ing Award recipients have begun publicly advocating for
safety research (Bengio et al., 2024).
Major venues
have recognized this work: NeurIPS lists safety, inter-
pretability, fairness, and privacy among its topical ar-
eas (NeurIPS, 2025b), hosts alignment and interpretability
workshops (NeurIPS, 2024b; Abrol et al., 2024), and re-
quires ethics review (NeurIPS, 2025a). Award committees
have followed, recognizing alignment and safety work at
NeurIPS 2023, NeurIPS 2024, and ICLR 2025 (NeurIPS,
5
Futures that benchmark-driven AI cannot see
2023; 2024a; ICLR, 2025). Leading laboratories have estab-
lished dedicated interpretability and alignment teams (Ope-
nAI, 2025; Anthropic, 2025; Google DeepMind, 2025).
Governments have created AI Safety Institutes across multi-
ple jurisdictions (UK Government, 2023; NIST, 2023; IAPS,
2025). These developments suggest a new set of concerns
is gaining attention. The 1956 Dartmouth proposal (Mc-
Carthy et al., 1956) asked: Can machines exhibit intelligent
behavior? The emerging question appears to extend this:
Can machines exhibit intelligent behavior such that they
are aligned, interpretable, and safe? We refer to this as the
such-that question.
4.2. Fundamentally Different from the Original
Question
The founding question of artificial intelligence—whether
machines can exhibit intelligent behavior— was deliberately
vague, but it admitted an epistemological bypass that made
rapid progress possible. Researchers could set aside the
philosophically fraught question of what intelligence is and
instead measure what systems do. Russell and Norvig’s
textbook codifies this move: the field studies rational agents
operating in task environments where success criteria can
be externally specified (Russell \& Norvig, 2021). Whether a
bird appears in a photograph admits ground truth. The rules
of Go are fixed before play begins. These are what we might
call world-side problems: the ground truth exists indepen-
dently of who evaluates the system, the task environment is
characterizable in advance, and closed-loop benchmarking
becomes a natural methodology. The contrast we draw is
not between objective and subjective questions—both kinds
admit precise formulations—but between problems whose
success criteria are fixed by the task and those that must be
constructed through deliberation.
Such-that clauses invert this structure. When we ask whether
a system behaves such that it remains aligned with human
values, interpretable to human understanding, or safe under
deployment, we are not asking whether it maps inputs to
outputs according to some function waiting to be discovered
in the world. We are asking whether its behavior accords
with standards that are contested, contextual, and constantly
renegotiated through social processes. The ground truth of
interpretability or alignment with society does not exist in
nature; it must be constructed through human deliberation.
Gabriel’s analysis of alignment identifies a “simple thesis”
that technologists sometimes assume, that one can solve
the technical problem first and specify normative criteria
afterward, and shows why it fails: normative and technical
dimensions are entangled from the start (Gabriel, 2020). The
such-that problem is human-side, not world-side, and the
difference is not one of degree. Here we focus specifically
on the normative subset of such-that clauses: alignment,
interpretability, safety, governance. Operational constraints
like energy efficiency or latency remain largely world-side
and benchmark-tractable. The distinction matters because
our argument targets the former, not the latter.
Q1 and Q2 are both design problems, but they differ in
whether the objective can be specified in advance (Simon,
1996). For Q1, success criteria are externally fixed and
agreed upon before the system is built. For Q2, they are
not: there is no single accepted definition of alignment,
interpretability, or safety. Gabriel documents multiple con-
flicting conceptions of alignment alone (Gabriel, 2020); the
interpretability literature lacks a settled account of what
interpretability requires (Lipton, 2018); and fairness criteria
have been shown to be mutually incompatible under reason-
able assumptions (Chouldechova, 2017). A field that cannot
agree on what it is trying to achieve cannot validate that
its benchmarks track the goal. Benchmark culture applied
to pre-paradigmatic problems does not accelerate progress
toward the goal; it accelerates progress toward whichever
operationalization happens to gain traction first. We have
built technology that outpaces the conceptual vocabulary
required to understand it. Hewitt, Geirhos, and Kim argue
that we cannot even describe AI behavior using existing
human language; interpretability itself is a communication
problem requiring the invention of new concepts (Hewitt
et al., 2025). The such-that problem demands not just better
engineering but something closer to a new science, one for
which the methods, and perhaps even the words, do not yet
exist.
Recent work from within the machine learning community
confirms this diagnosis. Fisher et al. argue that political
neutrality in AI is “theoretically impossible” because neu-
trality is inherently subjective. What seems neutral to one
perspective appears biased from another (Fisher et al., 2025).
Their framework draws explicitly on philosophy, political
science, and sociology, fields that have grappled with neu-
trality and bias for decades. Rahwan et al. call for a new
interdisciplinary field of “machine behavior” that would
study AI systems using methods from biology, psychology,
economics, and anthropology (Rahwan et al., 2019). These
are not peripheral voices; Fisher et al. received an oral pre-
sentation at ICML 2025, and Rahwan et al. published in
Nature. The message from inside the community is increas-
ingly clear: the such-that problem cannot be solved with the
tools that solved the capability problem.
4.3. The Case for Exaptation
If the such-that problem is genuinely different in kind, one
might expect the field’s response to differ accordingly. It has
not. Current attempts to benchmark alignment and safety
(refusal rates, bias metrics, RLHF/RLAIF scores, Constitu-
tional AI evaluations) represent genuine progress. Yet they
face a structural difficulty: we are optimizing proxies before
6
Futures that benchmark-driven AI cannot see
validating that they track the underlying goals. For capa-
bility research, decades of experience built confidence that
benchmark performance correlated with real-world progress.
For such-that problems, this correlation is precisely what
remains unestablished. The risk is not that safety cannot
be measured. The risk is that premature convergence on
current metrics forecloses discovery of what we actually
need to measure. The concepts required to measure it may
not yet exist.
History offers at least one suggestive precedent. The al-
gorithm now called backpropagation, foundational to mod-
ern deep learning, did not originate within artificial intel-
ligence research. Its mathematical core, the reverse mode
of automatic differentiation, first appeared in Linnainmaa’s
1970 master’s thesis on numerical error analysis, written
in Finnish with no reference to neural networks (Linnain-
maa, 1970). Related ideas appeared in optimal control
theory (Kelley, 1960). Werbos, whose 1974 PhD thesis
applied these techniques to neural network training, was
working not in computer science but in the behavioral sci-
ences, attempting to model political forecasting (Werbos,
1994). Rumelhart, Hinton, and Williams popularized the
method for machine learning over a decade later (Rumel-
hart et al., 1986). The lesson is not that “the solution was
waiting.” Considerable conceptual work was required to
recognize the relevance and bridge the gap. The deeper
lesson is that the field could not have specified in advance
what it needed, and therefore could not have benchmarked
toward it. Someone had to be working at the boundaries, in
contact with multiple traditions, for that bridge to be built at
all.
What might the analogous story look like for such-that prob-
lems? We do not know. Neither does anyone else. Bengio et
al. state explicitly that we should search for technical guar-
antees of safety if they exist, while acknowledging that we
do not know whether such guarantees exist (Bengio et al.,
2024). This is not ordinary uncertainty about parameters. It
is uncertainty about whether the problem admits the kind
of solution we are searching for. Under such conditions,
the exaptation tax becomes acute. Benchmark-driven se-
lection may be optimizing confidently in a subspace that
does not contain what we most need to find, while filtering
out exploratory work from which unexpected solutions have
historically emerged. The such-that problem may require
innovations in methods, but also in measurement regimes,
conceptual vocabulary, and institutional structure, domains
where benchmark-driven selection has less purchase.
Benchmarking itself was an institutional innovation that
enabled the capability era. The Common Task Framework
emerged from evaluation programs designed for specific
government needs, then was repurposed as a general se-
lection mechanism for ML (Koch \& Peterson, 2024). But
benchmarks harvested what prior exploration had planted.
Neural networks persisted through two AI winters despite a
selection environment that had largely moved on, sustained
by researchers who stayed at the boundary as connection-
ism became unfashionable and institutional support shifted
elsewhere. The deep learning boom required the selection
mechanism and something to select from. For capability
research, the accumulated exploration capital was sufficient.
For such-that problems, we cannot assume the same. The
concern is that if benchmark culture crowds out exploratory
work, the next era will have nothing to exapt.
5. Proposals
The proposals below outline institutional mechanisms for
restoring exaptive capacity without abandoning benchmark-
ing: pluralizing selection criteria, protecting exploratory
communities, and governing the selection rules themselves.
5.1. Plural Evaluation Regimes
Peer review in AI is nominally multi-criteria but decisions
are often made by collapsing those criteria into a single
ranking. When commonly used criteria such as novelty,
correctness, and empirical validation pull in different di-
rections, aggregation obscures what kind of contribution a
paper actually makes. An alternative is to run multiple scor-
ing regimes in parallel (e.g., one that highly rewards novelty
and another that rewards benchmark performance) and to
admit a portfolio that includes top papers from each regime.
Under such a system, a paper that introduces a new approach
but does not improve benchmark performance could still be
ranked highly in one regime and published. Venues such as
TMLR already instantiate part of this logic, evaluating sub-
missions on significance and correctness rather than novelty
thresholds or benchmark gains (Transactions on Machine
Learning Research, 2022). The such-that problem demands
more: venues that explicitly protect work whose goals are
themselves under construction.
A complementary mechanism is explicit diversification com-
bined with controlled randomness once a basic quality
threshold is met. Peer review in AI is noisy with substantial
variance across reviewers (Fogelholm et al., 2012; Aczel
et al., 2025). Under time pressure, reviewers often over-
weight easily checked signals such as benchmark perfor-
mance, which can distort judgments about criteria more
relevant to exaptive capacity (Alberts et al., 2014; Severin \&
Egger, 2021). When many reviewers apply the same short-
cut, the resulting bias becomes correlated noise, and partial
randomization among near-equal submissions can improve
selection (Kleinberg \& Raghavan, 2018; Emelianov et al.,
2020). Similar logic underlies modified lottery schemes
used in research funding (Adam et al., 2019; Liu et al.,
2020; Heyard et al., 2022) and structured hiring interven-
7
Futures that benchmark-driven AI cannot see
tions such as the Rooney Rule, which enforce diversity at
the consideration stage (Kleinberg \& Raghavan, 2018; Celis
et al., 2021). Stochastic selection is a standard tool for
maintaining diversity and avoiding premature convergence
in evolutionary computation and search theory (Lehman \&
Stanley, 2011; Emelianov et al., 2020).
Plural regimes should also treat disagreement as signal and
not failure. High variance in assessments often indicates
that a contribution challenges existing problem framings or
relies on premises not yet widely shared (Stanley \& Lehman,
2015). An idea that appears valuable to some experts but not
others is not a failed evaluation; it marks a frontier where
prevailing frameworks have not yet stabilized. Preserving
such disagreement allows ideas with latent potential to per-
sist long enough to realize it.
Notably, a similar pluralism has already emerged within AI
systems themselves (Sorensen et al., 2024). Contemporary
architectures increasingly rely on ensembles, routing, and
specialized components, with overall performance arising
from coordination rather than dominance by a single model.
The same community has developed novelty search, qual-
ity diversity algorithms, and uncertainty-aware exploration
methods precisely because pure objective optimization con-
verges too fast and too narrowly (Lehman \& Stanley, 2011).
Formal models of scientific communities confirm the trade-
off: more connected networks reach consensus faster, but
less connected ones are more likely to reach the correct
answer (Zollman, 2007). Research evaluation faces an anal-
ogous problem but applies the simplest possible method to
a task its own algorithms would approach very differently.
5.2. Protecting Exploratory Work and Enabling
Paradigm Turnover
Plural evaluation should be complemented by explicit mech-
anisms that protect communities with distinct standards of
progress. Ideas that challenge dominant assumptions rarely
perform well under prevailing metrics at early stages. With-
out protected environments in which alternative criteria are
legitimate, such ideas are eliminated before they can ma-
ture. Institutions should therefore support venues, funding
streams, and career signals that allow exploratory work to
persist even when misaligned with mainstream benchmarks.
In fast-growing fields, progress is often limited not by a lack
of new ideas, but by the slow displacement of dominant
frameworks that continue to perform well on existing met-
rics (Kuhn, 1997; Foster et al., 2015). Review and funding
processes should explicitly value work that stress-tests or
falsifies core assumptions, even when it does not offer an
immediate replacement or benchmark improvement.
Selection mechanisms must also account for the fact that in-
cremental advances within a paradigm compose cleanly
while shifts in paradigm typically require coordinated
changes across architectures, objectives, training methods,
and evaluation itself. Institutions should therefore tolerate
incomplete or uneven contributions that make sense only
as part of a broader reconfiguration, rather than penalizing
them for failing to deliver isolated gains.
Hochreiter and Schmidhuber published the LSTM in
1997 and watched the field move to support vector ma-
chines (Hochreiter \& Schmidhuber, 1997). The paper sat
slow to accumulate citations through a decade of benchmark-
driven dominance. It later powered the voice assistants that
reached production at scale and sustained the sequence mod-
eling tradition that made the Transformer possible. The
work this paper is trying to protect does not announce itself.
It looks, from the outside, like the LSTM looked in 2003:
real, rigorous, and invisible to the metrics that matter. If you
are working on a question the field cannot yet evaluate, you
are not in the wrong place. You are in the only place from
which certain kinds of progress can come.
5.3. Reflexive Governance of Evaluation Criteria
In addition to pluralizing selection and protecting ex-
ploratory communities, the research ecosystem must main-
tain the capacity to interrogate and revise its own selection
rules. Evaluation criteria in AI function as the main gover-
nance mechanisms: they define what counts as progress and
which failures are tolerable. Institutions should therefore
legitimize and support meta-level work on evaluation itself,
including empirical and conceptual analyses of benchmarks,
metrics, and problem framings. Critique of dominant selec-
tion rules should be treated as a core research activity with
dedicated venues and tracks explicitly designed for such
work such as the ICML position paper track. Initiatives
of this kind play an important role in making evaluative
assumptions explicit and open to revision. The stakes ex-
tend beyond the research community: benchmark culture
is now being encoded into regulatory frameworks for AI,
and evaluation assumptions that harden into legal standards
are an order of magnitude harder to revise than those that
remain community norms.
Finally, training and norms should equip researchers to rec-
ognize selection rules as contingent rather than fixed. This
includes mentoring around the practice of reviewing itself,
emphasizing judgment about ideas, assumptions, and po-
tential rather than reliance on surface metrics. Conferences
can reinforce these norms by recognizing reviewing quality
explicitly: ICML 2026 introduced Gold and Silver reviewer
designations based on review quality rather than throughput
alone. The community should go further: the most rigor-
ous, idea-engaged reviews, those that evaluate a contribu-
tion’s conceptual ambition and long-term potential, should
be citable scholarly contributions in their own right.
8
Futures that benchmark-driven AI cannot see
6. Alternative Views
The Success Objection. Benchmark culture has produced
AlphaFold, GPT-4, diffusion models, and the modern LLM
ecosystem. The field has never been more productive or
consequential. Where exactly is the problem? We do not
dispute this success, but we question its attribution. The
deep learning revolution required benchmarking as a se-
lection mechanism. It also required something to select
from. Neural networks persisted through two AI winters
despite a selection environment that had moved on: Rosen-
blatt’s perceptron in the 1950s, Hinton and Rumelhart in
the 1980s, Bengio and LeCun through the 1990s and 2000s
as the community shifted toward SVMs despite available
connectionist results. Benchmark-driven selection has har-
vested more than it has planted. For capability research,
the question of whether machines can exhibit intelligent
behavior, benchmark-driven selection has been extraordinar-
ily effective at harvesting accumulated exploration capital.
Our argument is forward-looking: the selection regime op-
timized for Q1 may not be optimized for Q2, the question
of whether machines can exhibit intelligent behavior such
that they are aligned, interpretable, and safe. These are dif-
ferent problems with different success criteria (Section 4.2).
Past performance on capability questions is not evidence of
future performance on such-that questions. The exaptation
tax is not visible in current achievements; it is visible—or
rather, invisible—in the paths not taken that might matter
for problems we are only beginning to face. We are not
arguing the system failed. We are arguing it succeeded at
one problem and now faces another.
The Invisibility Objection. If lost innovations cannot be
observed, the exaptation tax is unfalsifiable. Every exam-
ple cited—backpropagation, attention, RLHF—eventually
succeeded. We concede the specific counterfactual is un-
available, but this epistemic situation is familiar from con-
servation biology: one cannot identify which lost species
would have proven essential, yet the case for preserving bio-
diversity under uncertainty is widely accepted. Structural
narrowing is observable: thematic diversity in AI research
has stagnated, with increasing specialization toward compu-
tationally intensive methods (Klinger et al., 2020). We do
not claim benchmark culture is the sole cause—industrial
organization, compute costs, and convergence on effective
methods all contribute—but benchmark-driven selection
plausibly amplifies the narrowing by creating increasing re-
turns for a particular resource profile. Moreover, Lehman et
al. argue that ML systems face a “Knightian blindspot” when
confronted with qualitatively novel challenges (Lehman
et al., 2025); the same logic applies to evaluation systems
themselves. Under genuine uncertainty about which di-
rections will prove essential, reduced diversity is a cost in
expectation—even when specific losses remain unobserv-
able.
The Protection Objection. “Protected venues” could shel-
ter unproductive research from accountability. How do we
distinguish sleeping beauties from deservedly ignored work?
We cannot, and do not claim to. Protection means plural
evaluation, not absent evaluation: work must demonstrate
excellence under some criteria—rigor, conceptual clarity,
novelty—even if it does not improve benchmark perfor-
mance. Some protected work will fail; that is inherent to
exploration. The question is whether a diversified portfolio
outperforms a concentrated one when we cannot predict
which directions matter. We propose time-limited runway
and ex-post tracking, not permanent shelter.
The Industry Lab Objection. Elite labs already escape
benchmark pressure and produce transformative work. The
market has self-corrected; these proposals are unnecessary.
If the field’s exploratory capacity concentrates in organiza-
tions with patient capital, two concerns follow. First, which
questions get explored depends on which questions those or-
ganizations prioritize—a narrow governance structure for a
broadly consequential technology. The such-that questions
have public-goods characteristics that may be undersupplied
by organizations optimizing for competitive advantage. Sec-
ond, academic research risks losing relevance for the field’s
most important problems. Our proposals address whether
exploratory capacity should be rationed by resources or
distributed by design.
The Measurement Objection. Such-that problems need
more rigorous measurement, not less. RLHF succeeded
because it could be evaluated. The paper could justify aban-
doning rigor where rigor matters most. We agree that de-
veloping metrics for alignment, interpretability, and safety
is necessary. Our argument is not against measurement but
against treating metrics as dominant selection criteria be-
fore validating that they track our actual qualitative goals.
For capability research, experience gave confidence that
benchmark performance correlated with progress. For such-
that problems, this correlation is precisely what remains
unvalidated. The risk is mistaking measurability for valid-
ity—optimizing legible proxies while true objectives remain
underspecified. Metrics are tools; benchmark culture is a
selection regime, not one we endorse extending where the
metric-objective gap is not yet understood.
7. Discussion
Our argument is not that benchmarking is wrong but that
the current calibration is off: the field is over-exploiting a
selection regime optimized for Q1 at a moment when Q2
demands more exploration. But the asymmetry of error costs
matters. If we are wrong, the cost is some inefficiency. If we
are right, the field may be filtering out what it most needs
to find. Our proposals are not alternatives to benchmarking.
They are hedges against a future we cannot yet evaluate.
9
Futures that benchmark-driven AI cannot see
References
Abrol,
V.,
Boreiko,
V.,
Chaudhuri,
K.,
Di
Cas-
tro, D., Frost, N., Lakkaraju, H., Moshkovitz, M.,
Neuhof, B., Semenova, L., Srinivas, S., and Ya-
dav, C.
Interpretable AI: Past, present and fu-
ture.
NeurIPS 2024 Workshop, December 2024.
URL https://interpretable-ai-workshop.
github.io/. Workshop website; held Dec 15, 2024.
Accessed 2026-01-29.
Aczel, B., Barwich, A.-S., Diekman, A. B., Fishbach, A.,
Goldstone, R. L., Gomez, P., Gundersen, O. E., von Hip-
pel, P. T., Holcombe, A. O., Lewandowsky, S., et al. The
present and future of peer review: Ideas, interventions,
and evidence. Proceedings of the National Academy of
Sciences, 122(5):e2401232121, 2025.
Adam, D. et al. Science funders gamble on grant lotteries.
Nature, 575(7785):574–575, 2019.
Akhtar, M., Reuel, A., Soni, P., et al. When AI benchmarks
plateau: A systematic study of benchmark saturation.
arXiv preprint arXiv:2602.16763, 2026.
Alberts, B., Kirschner, M. W., Tilghman, S., and Varmus, H.
Rescuing us biomedical research from its systemic flaws.
Proceedings of the National Academy of Sciences, 111
(16):5773–5777, 2014.
Anthropic. Interpretability team. Organization website,
2025.
Anthwal, N., Joshi, L., and Tucker, A. S. Evolution of the
mammalian middle ear and jaw: adaptations and novel
structures. Journal of Anatomy, 222(1):147–160, 2013.
Arthur, W. B. Competing technologies, increasing returns,
and lock-in by historical events. The Economic Journal,
99(394):116–131, 1989.
Arthur, W. B. Increasing returns and the two worlds of busi-
ness. Harvard Business Review, 74(4):100–109, 1996.
Bahdanau, D., Cho, K., and Bengio, Y. Neural machine
translation by jointly learning to align and translate. arXiv
preprint arXiv:1409.0473, 2014.
Bechler-Speicher, M., Finkelshtein, B., Frasca, F., M¨uller,
L., T¨onshoff, J. M., Siraudin, A., Zaverkin, V., Bronstein,
M. M., Niepert, M., Perozzi, B., Galkin, M., and Morris,
C. Position: Graph learning will lose relevance due to
poor benchmarks. Proceedings of Machine Learning
Research, 267:81067–81089, 2025. URL https://
arxiv.org/abs/2502.14546.
Bengio, Y., Courville, A., and Vincent, P. Representation
learning: A review and new perspectives. IEEE Transac-
tions on Pattern Analysis and Machine Intelligence, 35
(8):1798–1828, 2013.
Bengio, Y. et al.
Managing ai risks in an era of rapid
progress. Science, 2024.
Bourdieu, P. The specificity of the scientific field and the
social conditions of the progress of reason. Social Science
Information, 14(6):19–47, 1975.
Celis, L. E., Hays, C., Mehrotra, A., and Vishnoi, N. K. The
effect of the rooney rule on implicit bias in the long term.
In Proceedings of the 2021 ACM conference on fairness,
accountability, and transparency, pp. 678–689, 2021.
Chouldechova, A. Fair prediction with disparate impact: A
study of bias in recidivism prediction instruments. Big
Data, 5(2):153–163, 2017.
Donoho, D. Data science at the singularity. Harvard Data
Science Review, 6(1), 2024.
doi: 10.1162/99608f92.
b91339ef.
Emelianov, V., Gast, N., Gummadi, K. P., and Loiseau, P.
On fair selection in the presence of implicit variance. In
Proceedings of the 21st ACM Conference on Economics
and Computation, pp. 649–675, 2020.
Everingham, M., Van Gool, L., Williams, C. K., Winn, J.,
and Zisserman, A. The PASCAL Visual Object Classes
(VOC) challenge. International Journal of Computer
Vision, 88(2):303–338, 2010.
Ferreira, M. R., Reisz, N., Schueller, W., Servedio, V. D. P.,
Thurner, S., and Loreto, V. Quantifying exaptation in
scientific evolution. arXiv preprint arXiv:2002.08144,
2020.
Fisher, J., Appel, R. E., Park, C. Y., Potter, Y., Jiang, L.,
Sorensen, T., Feng, S., Tsvetkov, Y., Roberts, M. E.,
Pan, J., Song, D., and Choi, Y. Political neutrality in AI
is impossible—but here is how to approximate it. Pro-
ceedings of the International Conference on Machine
Learning (ICML), 2025. Oral presentation.
Fogelholm, M., Leppinen, S., Auvinen, A., Raitanen, J.,
Nuutinen, A., and V¨a¨an¨anen, K. Panel discussion does
not improve reliability of peer review for medical research
grant proposals. Journal of clinical epidemiology, 65(1):
47–52, 2012.
Forcada, M. L. and ˜Neco, R. P.
Recursive hetero-
associative memories for translation. In International
Work-Conference on Artificial and Natural Neural Net-
works (IWANN), pp. 453–462. Springer, 1997.
Foster, J. G., Rzhetsky, A., and Evans, J. A. Tradition and
innovation in scientists’ research strategies. American
sociological review, 80(5):875–908, 2015.
Gabriel, I. Artificial intelligence, values, and alignment.
Minds and Machines, 30(3):411–437, 2020.
10
Futures that benchmark-driven AI cannot see
Garfield, E. Premature discovery or delayed recognition—
why? Current Contents, 21:5–10, 1980.
Goodfellow, I., Pouget-Abadie, J., Mirza, M., Xu, B.,
Warde-Farley, D., Ozair, S., Courville, A., and Bengio, Y.
Generative adversarial nets. In Advances in Neural Infor-
mation Processing Systems, volume 27, pp. 2672–2680,
2014.
Google DeepMind. Interpretability team. Organization
website, 2025.
Gould, S. J. and Vrba, E. S. Exaptation—a missing term in
the science of form. Paleobiology, 8(1):4–15, 1982.
Hao, Q., Xu, F., Li, Y., et al. Artificial intelligence tools
expand scientists’ impact but contract science’s focus. Na-
ture, 649:1237, 2026. URL https://www.nature.
com/articles/s41586-025-09922-y.
Harman, D. K.
The first text REtrieval conference
(TREC-1).
Technical Report Special Publication
500-207, National Institute of Standards and Technology,
Gaithersburg, MD, USA, May 1993.
URL https:
//nvlpubs.nist.gov/nistpubs/Legacy/
SP/nistspecialpublication500-207.pdf.
Hendrycks, D. and Dietterich, T. Benchmarking neural
network robustness to common corruptions and perturba-
tions. In International Conference on Learning Represen-
tations, 2019.
Hestness, J., Narang, S., Ardalani, N., Diamos, G., Jun,
H., Kianinejad, H., Patwary, M., Yang, Y., and Zhou, Y.
Deep learning scaling is predictable, empirically. arXiv
preprint arXiv:1712.00409, 2017.
Hewitt, J., Geirhos, R., and Kim, B.
We can’t under-
stand AI using our existing vocabulary. arXiv preprint
arXiv:2502.07586, 2025.
Heyard, R., Ott, M., Salanti, G., and Egger, M. Rethinking
the funding line at the swiss national science foundation:
Bayesian ranking and lottery. Statistics and Public Policy,
9(1):110–121, 2022.
Ho, J., Jain, A., and Abbeel, P. Denoising diffusion proba-
bilistic models. Neural Information Processing Systems
(NeurIPS), 33:6840–6851, 2020.
Hochreiter, S. and Schmidhuber, J. Long short-term memory.
Neural Computation, 9(8):1735–1780, 1997.
Hooker, S. The hardware lottery. CoRR, abs/2009.06489,
2020.
URL https://arxiv.org/abs/2009.
06489.
IAPS. International ai safety network. Organization website,
2025.
ICLR.
Iclr
2025
awards.
https://blog.iclr.cc/2025/04/22/announcing-the-
outstanding-paper-awards-at-iclr-2025/, 2025.
Jelinek, F. Statistical Methods for Speech Recognition. MIT
Press, Cambridge, MA, USA, 1997. ISBN 0262100665.
Jumper, J., Evans, R., Pritzel, A., Green, T., Figurnov, M.,
Ronneberger, O., Tunyasuvunakool, K., Bates, R., ˇZ´ıdek,
A., Potapenko, A., et al. Highly accurate protein structure
prediction with AlphaFold. Nature, 596(7873):583–589,
2021.
Kelley, H. J. Gradient theory of optimal flight paths. ARS
Journal, 30(10):947–954, 1960.
Kleinberg, J. and Raghavan, M.
Selection problems
in the presence of implicit bias.
arXiv preprint
arXiv:1801.03533, 2018.
Klinger, J., Mateos-Garcia, J., and Stathoulopoulos,
K.
A narrowing of AI research?
arXiv preprint
arXiv:2009.10385, 2020.
URL https://arxiv.
org/abs/2009.10385. Analysis of 110,000 AI pa-
pers from arXiv showing thematic diversity has stagnated;
private sector research less diverse than academia.
Knight, F. H.
Risk, Uncertainty and Profit. Houghton
Mifflin, Boston, MA, 1921.
Koch, B. J. and Peterson, D. From protoscience to epis-
temic monoculture: How benchmarking set the stage for
the deep learning revolution. https://arxiv.org/
abs/2404.06647, 2024. arXiv:2404.06647.
Krizhevsky, A., Sutskever, I., and Hinton, G. E. ImageNet
classification with deep convolutional neural networks.
In Advances in Neural Information Processing Systems,
volume 25, 2012.
Kuhn, T. S. The Structure of Scientific Revolutions. Univer-
sity of Chicago Press, 1962.
Kuhn, T. S. The structure of scientific revolutions, volume
962. University of Chicago press, 1997.
Langley, P. Crafting papers on machine learning. In Langley,
P. (ed.), Proceedings of the 17th International Conference
on Machine Learning (ICML 2002), pp. 1207–1216, Stan-
ford, CA, 2002. Morgan Kaufmann.
LeCun, Y., Bengio, Y., and Hinton, G. Deep learning. Na-
ture, 521(7553):436–444, 2015.
Lehman, J. and Stanley, K. O. Abandoning objectives: Evo-
lution through the search for novelty alone. Evolutionary
computation, 19(2):189–223, 2011.
11
Futures that benchmark-driven AI cannot see
Lehman, J., Meyerson, E., El-Gaaly, T., Stanley, K. O., and
Ziyaee, T. Evolution and the Knightian blindspot of ma-
chine learning. arXiv preprint arXiv:2
\end{document}
